\chapter*{第一册能力清单}
\addcontentsline{toc}{chapter}{第一册能力清单}

读完第一册并完成实验后，能力不应停在“知道编译流程”或“看过汇编”。最低出口是能把一个小型 C++ kernel 从源码追到正在运行的进程，并能说明每层证据的边界。

\section*{能实现}

\begin{itemize}
  \item 写一个带 reference、candidate、checksum、CSV 和环境字段的小型 benchmark。
  \item 写一个 \code{TensorViewF32} 级别的对象布局实验，输出 \code{sizeof}、\code{alignof}、成员偏移和 buffer 地址。
  \item 写一个 C++ 调用 C ABI 或手写汇编 wrapper 的最小程序，并保持 System V AMD64 ABI 正确。
  \item 写 sanitizer 和 fuzz-style shape regression，覆盖越界、坏 stride、空输入、tail 和溢出形状。
\end{itemize}

\section*{能诊断}

\begin{itemize}
  \item 用 \code{readelf}、\code{nm}、\code{objdump}、GDB 和 \filepath{/proc/<pid>/maps} 判断运行的是哪份二进制。
  \item 区分源码函数、目标文件符号、ELF 动态链接、PLT/GOT、loader 和进程地址空间。
  \item 判断一个 benchmark 是否被优化器删除、是否混入启动成本、I/O、首次缺页或动态链接成本。
  \item 解释 \code{optimized out}、frame pointer 省略、DWARF CFI、异常展开和 backtrace 质量的关系。
\end{itemize}

\section*{能解释}

\begin{itemize}
  \item 从 C++ 语义追到 AST/IR、SSA、load/store、phi、DCE、LICM、inline 和 vectorization remark。
  \item 从函数签名推导参数寄存器、返回寄存器、栈对齐、caller/callee saved 和 red zone。
  \item 从结构体定义推导 padding、alignment、object lifetime、strict aliasing、pointer provenance 和 UB 风险。
  \item 把 syscall、vDSO、minor/major page fault、COW、signal 和用户态/内核态边界分开。
\end{itemize}

\section*{不能夸大的部分}

第一册只建立单机程序到机器执行的基础。若没有第二册的 cache/TLB/NUMA/SIMD/PMU 和第三册的模型/KV/量化/runtime 训练，不能把第一册能力直接等同于完整性能工程或 AI Infra 能力。
